\frontchapter{VIII. Sources and Illustration Design}
\section*{1. Primary Materials Required in Each Part}
\begin{itemize}
\item A photograph of an object or archaeological site.
\item A contemporary map or later historical reconstruction.
\item An extract from a law, contract, or administrative record.
\item An extract from a poem, story, letter, or diary.
\item A voice from a marginalized group.
\item A statistical chart or data on prices, population, or trade.
\item Two conflicting accounts of the same event.
\end{itemize}
Use only necessary short quotations from modern works still under copyright, favoring summaries, close readings of brief passages, and authorized material. Classical texts should also identify their translators and editions.

\section*{2. A Consistent Visual Grammar}
\begin{itemize}
\item Use timeline colors to distinguish regions, never levels of civilizational achievement.
\item Show political borders, ecological regions, and exchange routes together on maps.
\item On profile pages, list what a person proposed, what evidence they offered, and how later thinkers criticized them.
\item Use different arrows in argument diagrams for support, rebuttal, qualification, and counterexamples.
\item Present wording, narrative viewpoint, structure, and historical context in separate columns on close-reading pages.
\item Compare questions, concepts, and practices in tables of religions and philosophies, without scoring superiority.
\end{itemize}
